\documentclass[11pt,letterpaper]{article}
\usepackage[margin=0.75in]{geometry}
\usepackage{booktabs}
\usepackage{array}
\usepackage{xcolor}
\usepackage{colortbl}
\usepackage{caption}
\usepackage{hyperref}
\usepackage{fancyhdr}
\usepackage{graphicx}

\definecolor{toprow}{HTML}{E8F5E9}
\definecolor{midrow}{HTML}{FFF8E1}
\definecolor{botrow}{HTML}{FFEBEE}
\definecolor{hdrblue}{HTML}{1565C0}
\definecolor{notegray}{HTML}{757575}
\definecolor{warnrow}{HTML}{FFF3E0}

\pagestyle{fancy}
\fancyhf{}
\rhead{\textcolor{notegray}{\small March 11, 2026}}
\lhead{\textcolor{notegray}{\small SLM Agentic Benchmarking}}
\cfoot{\thepage}

\newcolumntype{R}[1]{>{\raggedleft\arraybackslash}p{#1}}
\newcolumntype{C}[1]{>{\centering\arraybackslash}p{#1}}

\begin{document}

\begin{center}
{\LARGE \textbf{SLM Agentic Benchmark Results}}\\[4pt]
{\large March 11, 2026}\\[6pt]
{\small\textcolor{notegray}{Benchmark: BIG-bench Lite 24-task Suite (BBL24)}}\\[2pt]
{\small\textcolor{notegray}{Models: 8 Ollama SLMs (0.6B--20B) $\cdot$ 4 Architectures (One-shot, Sequential, Concurrent, Group Chat)}}\\[2pt]
{\small\textcolor{notegray}{Infrastructure: 8$\times$ Modal L4 GPU workers $\cdot$ 10 examples/task $\cdot$ 308 questions/combo}}
\end{center}

\vspace{0.3cm}

%% =======================================================================
\section*{Methods}
%% =======================================================================

\textbf{BIG-bench Lite 24 (BBL24).} A curated 24-task subset of Google's BIG-bench benchmark
(\texttt{tasksource/bigbench} on HuggingFace). Tasks cover diverse reasoning capabilities:
multiple-choice commonsense (BBQ, conceptual combinations, emoji movies, Hindu knowledge,
known unknowns, misconceptions Russian, novel concepts, play dialog, strange stories, StrategyQA,
symbol interpretation, VitaminC fact verification, WinoWhy), structured inference (formal
fallacies, logic grid puzzles, logical deduction), code understanding (auto-debugging, code line
description), linguistic analysis (conlang translation, language identification, linguistics
puzzles), math/operators (operators), and sequence generation (repeat-copy logic, ParsiNLU
reading comprehension). Scoring is auto-selected per task: multiple-choice letter match, exact
string match, or BLEU.

Task sizes vary: three tasks (\texttt{auto\_debugging}, \texttt{novel\_concepts},
\texttt{repeat\_copy\_logic}) have fewer than 10 examples in the dataset and contribute their
full available count (24--34 examples); the remaining 21 tasks use exactly 10 examples each.
Total questions per run: 308.

\textbf{Evaluation metric.} \emph{Task-equal-weighted accuracy} (\texttt{weighted\_accuracy}):
each task's per-example accuracy is averaged, then those 24 per-task averages are averaged
equally, giving equal weight to every task regardless of example count. This avoids
\texttt{auto\_debugging} (34 examples) dominating the aggregate.

\textbf{Models.} Eight Ollama models spanning 0.6B to 20B parameters were evaluated. Models
range from dense instruct-tuned transformers (Gemma3-1B, Gemma3-4B) to efficient mixture-of-
experts architectures (Gemma3n-E2B, Gemma3n-E4B), a specialized distilled reasoning model
(DASD-4B), a mini reasoning chain model (Phi4-Mini-Reasoning), a compact next-token predictor
(Qwen3-0.6B), and a 20B open-source model (GPT-OSS-20B).

\textbf{Architectures.} Four agent architectures were evaluated:
\begin{itemize}
  \item \textbf{One-shot}: Direct single-call OllamaAgent. Task $\to$ Ollama REST API $\to$ post-process $\to$ response. No agentic overhead. True non-agentic baseline.
  \item \textbf{Sequential}: Three-stage CrewAI pipeline (Analyzer $\to$ Evaluator $\to$ Responder). Each stage makes one LLM call; output feeds the next stage. 60-second per-question timeout.
  \item \textbf{Concurrent}: Four CrewAI agents (Analyst + Researcher + Critic $\to$ Synthesizer). The first three work independently, then a Synthesizer combines their outputs. 60-second timeout.
  \item \textbf{Group Chat}: Discussion-style pipeline (Proposer $\to$ Critic $\to$ Advisor $\to$ Moderator). 60-second timeout.
\end{itemize}

\textbf{Infrastructure.} Sweep ran on 8 parallel Modal L4 (24\,GB VRAM) GPU workers, one model
per worker. Each worker started a local Ollama server, pulled model weights into a persistent
volume (\texttt{slm-ollama-model-cache}), and ran the sweep sequentially across architectures.
Results were committed to \texttt{slm-bigbench-results} after each architecture completed.
The sweep was interrupted by a Modal billing cycle spending limit after approximately 8 hours,
before concurrent and group-chat architectures were reached for most models. As a result,
this report covers complete one-shot data for 7 models and complete sequential data for 6 models.

\textbf{Timeout behavior.} CrewAI's \texttt{kickoff\_with\_timeout} enforces a 60-second wall-
clock limit per question. On timeout, the framework calls \texttt{executor.shutdown(wait=False)},
which returns immediately but leaves worker threads running. For small models that reliably
timeout, ghost threads accumulate and compete for the Ollama server, causing average latency to
converge to exactly 90 seconds (the combined timeout + thread overhead). Per-question \texttt
{timed\_out: True} is recorded in \texttt{results.jsonl}; timed-out questions receive an empty
response which scores as incorrect.

\newpage

%% =======================================================================
\section*{1\quad One-Shot Results (7 Models Complete)}
%% =======================================================================

\noindent\textbf{TL;DR:} Efficient MoE architectures punch well above their nominal parameter
count. Gemma3n-E2B (2B active) leads at 51.5\%, beating the dense 4B models. Gemma3-4B reaches
47.1\%, the top dense model. Most sub-2B models cluster in the 25--32\% range, roughly double
random-chance on these 4-class MC tasks ($\approx$25\%).

\vspace{0.2cm}

\begin{table}[h!]
\centering
\caption*{\textbf{Table 1.} One-shot BBL24 accuracy, $n=308$ per model (except Falcon-H1-90M: $n=518$, 20 ex/task).
  Sorted by weighted accuracy. Green = top tier ($\geq$45\%), yellow = mid-tier (30--45\%), red = below chance ($<$25\%).}
\begin{tabular}{l C{0.8cm} R{1.5cm} R{1.5cm} R{1.5cm} R{2.0cm}}
\toprule
\textbf{Model} & \textbf{Params} & \textbf{Wtd Acc (\%)} & \textbf{Raw Acc (\%)} & \textbf{Avg Lat (s)} & \textbf{Total Time} \\
\midrule
\rowcolor{toprow} Gemma3n-E2B     & 2B MoE & 51.5 & 45.8 & 3.3  & 0.3\,h \\
\rowcolor{toprow} Gemma3-4B       & 4B     & 47.1 & 42.2 & 2.6  & 0.2\,h \\
\rowcolor{midrow} DASD-4B         & 4B     & 31.9 & 28.9 & 8.0  & 0.7\,h \\
\rowcolor{midrow} Phi4-Mini-Rea.  & 4B     & 31.8 & 28.6 & 7.9  & 0.7\,h \\
\rowcolor{midrow} Qwen3-0.6B      & 0.6B   & 31.3 & 28.6 & 3.9  & 0.3\,h \\
\rowcolor{midrow} Gemma3-1B       & 1B     & 27.3 & 25.0 & 2.1  & 0.2\,h \\
\rowcolor{botrow} Falcon-H1-90M   & 90M    & 24.9 & 23.4 & 2.2  & 0.3\,h \\
\bottomrule
\end{tabular}
\end{table}

\noindent\textbf{Key observations (one-shot):}
\begin{itemize}
  \item \textbf{MoE efficiency}: Gemma3n-E2B with only $\sim$2B active parameters (out of a
    larger total) outperforms all dense models including Gemma3-4B. This confirms that
    mixture-of-experts routing provides disproportionate reasoning gains per active compute.
  \item \textbf{Near-random sub-1B models}: Falcon-H1-90M (24.9\%) and Gemma3-1B (27.3\%) are
    barely above the 25\% random baseline for 4-class multiple-choice tasks. These models
    reliably extract letter choices but lack deeper reasoning.
  \item \textbf{Thinking overhead}: DASD-4B and Phi4-Mini-Reasoning score nearly identically
    (31.9\% vs 31.8\%) despite different architectures, with 4$\times$ higher latency than
    Gemma3-4B. Their reasoning chain overhead does not translate to higher accuracy on these tasks.
  \item \textbf{Qwen3-0.6B anomaly}: Despite being the smallest model (0.6B), Qwen3 ties the
    much larger DASD-4B on raw accuracy (28.6\%). Its efficient tokenization and strong
    instruction following compensate for smaller capacity.
\end{itemize}

\newpage

%% =======================================================================
\section*{2\quad Sequential Architecture: Capability Cliff}
%% =======================================================================

\noindent\textbf{TL;DR:} Sequential agentic architecture (Analyzer $\to$ Evaluator $\to$ Responder
via CrewAI) catastrophically degrades performance on every model below a capability threshold.
Models at or below 4B dense parameters time out on 92--100\% of questions, scoring at or near 0\%.
Only Gemma3n-E4B, an efficient 4B MoE model, crosses the threshold with 30\% timeouts and retains
meaningful accuracy (39.2\%).

\vspace{0.2cm}

\begin{table}[h!]
\centering
\caption*{\textbf{Table 2.} Sequential architecture BBL24 accuracy, $n=308$ per model.
  Timeout = fraction of questions where CrewAI exceeded the 60-second limit.
  Red = failed ($<$5\% accuracy). Green = functional ($\geq$30\%).}
\begin{tabular}{l C{0.8cm} R{1.5cm} R{1.5cm} R{1.5cm} R{1.8cm}}
\toprule
\textbf{Model} & \textbf{Params} & \textbf{Wtd Acc (\%)} & \textbf{Timeout (\%)} & \textbf{Avg Lat (s)} & \textbf{Total Time} \\
\midrule
\rowcolor{toprow} Gemma3n-E4B     & 4B MoE & 39.2 & 29.9 & 50.9 & 4.4\,h \\
\midrule
\rowcolor{botrow} Gemma3-1B       & 1B     &  3.5 & 92.5 & 91.0 & 7.8\,h \\
\rowcolor{botrow} Gemma3-4B       & 4B     &  0.5 & 98.4 & 89.1 & 7.6\,h \\
\rowcolor{botrow} Gemma3n-E2B     & 2B MoE &  0.5 & 96.8 & 88.4 & 7.6\,h \\
\rowcolor{botrow} Qwen3-0.6B      & 0.6B   &  0.0 & 100.0 & 90.0 & 7.7\,h \\
\rowcolor{botrow} GPT-OSS-20B     & 20B    &  0.0 & 100.0 & 90.0 & 7.7\,h \\
\bottomrule
\end{tabular}
\end{table}

\noindent\textbf{Key observations (sequential):}
\begin{itemize}
  \item \textbf{Hard capability cliff}: There is a binary divide between models that can
    sustain the CrewAI ReAct format and those that cannot. Models below the threshold fail
    entirely regardless of size --- the 20B GPT-OSS-20B times out on 100\% of questions,
    same as the 0.6B Qwen3.
  \item \textbf{Root cause}: CrewAI's \texttt{Process.sequential} uses a ReAct-style
    \texttt{Thought:/Action:/Observation:} format. Small models reliably produce prose answers
    but fail to emit the required \texttt{Action:} token within the timeout. Each failed
    step triggers a full 60-second wait before the framework gives up, driving average
    latency to 88--91 seconds per question.
  \item \textbf{Gemma3n-E4B exception}: The efficient MoE 4B model achieves only 30\%
    timeouts and reaches 39.2\% accuracy --- actually \textit{higher} than its one-shot result
    would predict. This suggests that when the agentic pipeline functions, the multi-stage
    reasoning provides genuine benefit.
  \item \textbf{GPT-OSS-20B paradox}: A 20B parameter model fails 100\% of sequential
    questions. This indicates that raw parameter count is not sufficient; what matters is
    the ability to follow structured output formats (JSON/ReAct) reliably within time
    constraints. The 20B quantized GGUF model likely produces slower tokens per second
    than smaller dense models, hitting the timeout before completing any ReAct step.
  \item \textbf{Architecture cost}: For non-functional models, sequential adds 7.6--7.8 hours
    of compute cost to produce \emph{worse} results than one-shot. The per-question cost
    is $\approx 30\times$ higher (90s vs 3s) with accuracy near zero.
\end{itemize}

\newpage

%% =======================================================================
\section*{3\quad Architecture Comparison (One-Shot vs.\ Sequential)}
%% =======================================================================

\noindent For models with both one-shot and sequential results, the relative accuracy impact:

\vspace{0.2cm}

\begin{table}[h!]
\centering
\caption*{\textbf{Table 3.} One-shot vs.\ sequential accuracy delta per model.
  $\Delta$ = Sequential $-$ One-shot. Positive = sequential helps; negative = hurts.}
\begin{tabular}{l C{0.8cm} R{1.4cm} R{1.4cm} R{1.4cm} R{2.5cm}}
\toprule
\textbf{Model} & \textbf{Params} & \textbf{1-Shot (\%)} & \textbf{Seq (\%)} & \textbf{$\Delta$ (pp)} & \textbf{Seq Timeout (\%)} \\
\midrule
\rowcolor{toprow} Gemma3n-E4B  & 4B MoE & ---  & 39.2 & --- & 29.9 \\
\rowcolor{midrow} Gemma3-4B    & 4B     & 47.1 &  0.5 & $-$46.6 & 98.4 \\
\rowcolor{midrow} Gemma3n-E2B  & 2B MoE & 51.5 &  0.5 & $-$51.0 & 96.8 \\
\rowcolor{midrow} Gemma3-1B    & 1B     & 27.3 &  3.5 & $-$23.8 & 92.5 \\
\rowcolor{botrow} Qwen3-0.6B   & 0.6B   & 31.3 &  0.0 & $-$31.3 & 100.0 \\
\rowcolor{botrow} GPT-OSS-20B  & 20B    & ---  &  0.0 & --- & 100.0 \\
\bottomrule
\end{tabular}
\end{table}

\noindent\textit{Note}: Gemma3n-E4B one-shot and GPT-OSS-20B one-shot were not completed before the
billing limit was reached.

\vspace{0.4cm}

\noindent\textbf{Latency overhead breakdown:}

\begin{table}[h!]
\centering
\caption*{\textbf{Table 4.} Latency cost per question: one-shot vs.\ sequential. Sequential
  overhead = avg sequential latency $\div$ avg one-shot latency.}
\begin{tabular}{l R{2.0cm} R{2.0cm} R{2.0cm}}
\toprule
\textbf{Model} & \textbf{1-Shot Lat (s)} & \textbf{Seq Lat (s)} & \textbf{Overhead ($\times$)} \\
\midrule
Gemma3-1B   &  2.1 & 91.0 & 43$\times$ \\
Gemma3-4B   &  2.6 & 89.1 & 34$\times$ \\
Gemma3n-E2B &  3.3 & 88.4 & 27$\times$ \\
Qwen3-0.6B  &  3.9 & 90.0 & 23$\times$ \\
Gemma3n-E4B &  ---  & 50.9 & --- \\
GPT-OSS-20B &  ---  & 90.0 & --- \\
\bottomrule
\end{tabular}
\end{table}

\noindent The sequential overhead for failed models is \emph{entirely} timeout waste. At 90 seconds
per question and 308 questions, each failed sequential run consumes $\approx$7.7 hours of GPU time
to produce zero correct answers.

\newpage

%% =======================================================================
\section*{4\quad Cross-Benchmark Skill Correlation}
%% =======================================================================

\noindent The following scatter plots show BIG-bench Lite accuracy (y-axis) vs.\ scores on
five prior skill benchmarks (x-axis): Planning, Criticality, Recall, Summarization, and
Instruction Following. Each point represents one (model, architecture) pair; color encodes
model, marker shape encodes architecture. \textbf{Hollow markers} indicate ``failed'' runs
(BBL24 $<$5\% accuracy OR skill score = 0), where the model produced no meaningful signal.
Ellipses show approximate 95\% confidence intervals (1.96$\times$SE in each axis).

All axes are fixed to 0--100\% range. Skill scores were multiplied by 100 for consistency.

\vspace{0.3cm}

\begin{figure}[h!]
\centering
\includegraphics[width=0.82\textwidth]{../../plots/bigbench_skill/skill_vs_bigbench_planning.png}
\caption*{\textbf{Figure 1.} BBL24 vs.\ Planning skill. Hollow markers = failed runs (BBL $<$5\% or planning = 0).}
\end{figure}

\begin{figure}[h!]
\centering
\includegraphics[width=0.82\textwidth]{../../plots/bigbench_skill/skill_vs_bigbench_criticality.png}
\caption*{\textbf{Figure 2.} BBL24 vs.\ Criticality skill.}
\end{figure}

\newpage

\begin{figure}[h!]
\centering
\includegraphics[width=0.82\textwidth]{../../plots/bigbench_skill/skill_vs_bigbench_recall.png}
\caption*{\textbf{Figure 3.} BBL24 vs.\ Recall skill.}
\end{figure}

\begin{figure}[h!]
\centering
\includegraphics[width=0.82\textwidth]{../../plots/bigbench_skill/skill_vs_bigbench_summarization.png}
\caption*{\textbf{Figure 4.} BBL24 vs.\ Summarization skill.}
\end{figure}

\begin{figure}[h!]
\centering
\includegraphics[width=0.82\textwidth]{../../plots/bigbench_skill/skill_vs_bigbench_instruction_following.png}
\caption*{\textbf{Figure 5.} BBL24 vs.\ Instruction Following skill. X-axis fixed 0--100 (percentage).
  Several failed sequential runs cluster at (0, 0) and are shown as hollow markers.}
\end{figure}

\newpage

%% =======================================================================
\section*{5\quad Summary and Key Findings}
%% =======================================================================

\subsection*{5.1 One-Shot Rankings}

Within the one-shot architecture, model capability follows parameter efficiency more closely than
raw count:

\begin{enumerate}
  \item \textbf{Gemma3n-E2B (51.5\%)} --- Best overall. MoE routing concentrates capacity on
    the active reasoning path, enabling 2B-active performance comparable to dense 4B models.
  \item \textbf{Gemma3-4B (47.1\%)} --- Best dense model. Strong across structured tasks
    (code, conceptual combinations, emoji movies at 100\%; operators at 90\%).
  \item \textbf{DASD-4B (31.9\%)} and \textbf{Phi4-Mini-Reasoning (31.8\%)} --- Thinking models
    add latency (8s vs 2.6s) without accuracy gains on these tasks.
  \item \textbf{Qwen3-0.6B (31.3\%)} --- Remarkable for 0.6B: ties the 4B thinking models,
    suggesting strong instruction adherence compensates for capacity limits.
  \item \textbf{Gemma3-1B (27.3\%)} and \textbf{Falcon-H1-90M (24.9\%)} --- At or near random
    baseline; reasoning tasks require minimum capacity not present at $<$1B parameters.
\end{enumerate}

\subsection*{5.2 The Agentic Overhead Problem}

Sequential (and by extension concurrent and group-chat) architectures impose a \textbf{hard
capability requirement} on models: they must reliably produce structured \texttt{Action:}
tokens under time pressure. For models that cannot:

\begin{itemize}
  \item \textbf{100\% timeout rate} $\Rightarrow$ 0\% accuracy, regardless of model size.
  \item \textbf{Per-question cost} jumps from 2--8 seconds (one-shot) to 88--91 seconds
    (sequential), a 23--43$\times$ overhead with \emph{zero return}.
  \item \textbf{Ghost thread accumulation} further degrades performance mid-run: each timed-out
    question leaves a running background thread that consumes Ollama connections, creating a
    feedback loop that worsens as the run progresses.
\end{itemize}

\subsection*{5.3 Dense vs.\ Efficient 4B: A Case Study}

The Gemma3-4B vs.\ Gemma3n-E4B comparison is the most instructive finding of this sweep:

\begin{center}
\begin{tabular}{l r r r r}
\toprule
Model & Params & 1-Shot Acc & Seq Acc & Seq Timeout \\
\midrule
Gemma3-4B   & 4B dense & 47.1\% & 0.5\%  & 98.4\% \\
Gemma3n-E4B & 4B MoE   & ---    & 39.2\% & 29.9\% \\
\bottomrule
\end{tabular}
\end{center}

\noindent Despite the same nominal 4B parameter count, the efficient MoE architecture produces
\textbf{$>$78$\times$ more correct answers} in the sequential setting. The MoE routing allows
the model to dedicate sparse expert capacity to format adherence, whereas the dense model
exhausts its capacity on content generation and fails to complete the ReAct formatting loop.

\subsection*{5.4 Implications for Model Selection}

\begin{itemize}
  \item \textbf{If using one-shot}: Choose the highest-capability model that fits latency
    constraints. Gemma3n-E2B is optimal below 3B active parameters.
  \item \textbf{If using agentic architectures}: Only models that cross the ReAct format
    threshold are viable. Currently confirmed: Gemma3n-E4B. Likely viable (untested here):
    larger dense models ($\geq$7B) and other MoE architectures with fast token generation.
  \item \textbf{Thinking models}: DASD-4B and Phi4-Mini-Reasoning's extended reasoning chains
    add latency without accuracy gains on BBL24 one-shot tasks. Their value may emerge on
    harder mathematical/logical tasks not covered in this benchmark.
  \item \textbf{Sub-1B models}: Useful only for latency-critical applications where near-random
    accuracy is acceptable. Not viable for any multi-step reasoning task.
\end{itemize}

\subsection*{5.5 Data Completeness Caveats}

Due to the Modal billing limit interruption, the following data is \textbf{missing} from this
report:
\begin{itemize}
  \item One-shot results for: Gemma3n-E4B, GPT-OSS-20B (incomplete runs)
  \item Concurrent architecture results: all models (not started before billing kill)
  \item Group-chat architecture results: all models (not started before billing kill)
  \item Sequential results for: DASD-4B, Falcon-H1-90M, Phi4-Mini-Reasoning (not reached)
\end{itemize}

A follow-up sweep with budget controls (per-worker timeouts, architecture priority ordering)
should complete the missing cells.

\end{document}
